% Cantilever beam bending — analytical derivation
% AIM Interactive Mechanics series. Compile with pdflatex.
\documentclass[10pt,a4paper]{article}
\usepackage[margin=2cm]{geometry}
\usepackage{amsmath}
\usepackage{amssymb}
\usepackage[T1]{fontenc}
\usepackage{lmodern}
\usepackage{microtype}
\usepackage{titlesec}
\titlespacing*{\section}{0pt}{9pt}{2pt}
\usepackage{xcolor}
\definecolor{iron}{HTML}{935636}
\definecolor{celdark}{HTML}{4B6C61}
\definecolor{celline}{HTML}{CDE0D8}
\usepackage[colorlinks=true,urlcolor=iron,linkcolor=black]{hyperref}
\setlength{\parindent}{0pt}
\setlength{\parskip}{4pt}
\AtBeginDocument{\setlength{\abovedisplayskip}{6pt}\setlength{\belowdisplayskip}{6pt}\setlength{\abovedisplayshortskip}{3pt}\setlength{\belowdisplayshortskip}{3pt}}
\pagestyle{empty}

\begin{document}

{\footnotesize\textcolor{celdark}{\textsc{aim interactive mechanics \,\textperiodcentered\, analytical derivation}}}

{\LARGE Cantilever Beam Bending}

\textcolor{celdark}{\small Companion to the interactive simulator at
\url{https://aimlab.uk/sims/cantilever-bending.html}}

\medskip
{\color{celline}\hrule height 1pt}

\section*{Setup and assumptions}
A prismatic cantilever of length $L$ and rectangular cross-section $w \times h$
(width $w$, height $h$) is clamped at $x = 0$ and carries a transverse point load
$F$ at the free end $x = L$. The material is linear elastic with Young's modulus
$E$. We measure the deflection $u(x)$ positive in the direction of the load and
adopt the Euler--Bernoulli assumptions: cross-sections remain plane and
perpendicular to the beam axis (shear deformation neglected), and deflections are
small ($\delta \lesssim 0.1\,L$).

\section*{Section property}
The second moment of area of the rectangle about its horizontal centroidal axis is
\begin{equation}
I \;=\; \int_A y^2 \,\mathrm{d}A \;=\; w \int_{-h/2}^{h/2} y^2 \,\mathrm{d}y
\;=\; \frac{w h^3}{12}.
\end{equation}

\section*{Governing equation}
Moment--curvature (from plane sections and Hooke's law) gives $M = EI\,u''$.
Equilibrium of a beam slice gives $\mathrm{d}M/\mathrm{d}x = V$ and
$\mathrm{d}V/\mathrm{d}x = -q$. Between loads $q = 0$, so
\begin{equation}
EI \,\frac{\mathrm{d}^4 u}{\mathrm{d}x^4} \;=\; 0 .
\end{equation}
Integrating four times introduces four constants:
\begin{equation}
u''' = C_1, \qquad
u'' = C_1 x + C_2, \qquad
u' = \tfrac{1}{2}C_1 x^2 + C_2 x + C_3, \qquad
u = \tfrac{1}{6}C_1 x^3 + \tfrac{1}{2}C_2 x^2 + C_3 x + C_4 .
\end{equation}

\section*{Boundary conditions}
The clamped end can neither deflect nor rotate; the free end carries no bending
moment, and its shear force equals the applied load:
\begin{equation}
u(0) = 0, \qquad u'(0) = 0, \qquad EI\,u''(L) = 0, \qquad EI\,u'''(L) = -F .
\end{equation}
The first two give $C_4 = C_3 = 0$; the last two give $C_1 = -F/EI$ and
$C_2 = FL/EI$.

\section*{Solution}
\begin{equation}
\boxed{\; u(x) \;=\; \frac{F x^2 \, (3L - x)}{6\,EI} \;}
\qquad\Rightarrow\qquad
\delta \;=\; u(L) \;=\; \frac{F L^3}{3\,EI},
\qquad
\theta_L \;=\; u'(L) \;=\; \frac{F L^2}{2\,EI}.
\end{equation}
The bending moment is largest at the wall, $M_{\max} = FL$, so the peak bending
stress (at the outer fibres, $y = \pm h/2$) is
\begin{equation}
\sigma_{\max} \;=\; \frac{M_{\max}\,(h/2)}{I} \;=\; \frac{6\,F L}{w h^2}.
\end{equation}
Deflection grows with $L^3$ and shrinks with $E$ and $h^3$: doubling the height
is eight times stiffer, doubling the width only twice.

\section*{Worked example (simulator defaults)}
$E = 2$~GPa, $L = 500$~mm, $w = 25$~mm, $h = 10$~mm, $F = 4$~N:
\begin{equation}
I = \frac{25 \times 10^3}{12} = 2083~\text{mm}^4, \qquad
\delta = \frac{4 \times 0.5^3}{3 \times (2\times10^{9}) \times (2.083\times10^{-9})}
\approx 40~\text{mm},
\end{equation}
with $\sigma_{\max} = 4.8$~MPa. Set the simulator sliders to these values and
compare.

\vfill
{\color{celline}\hrule height 0.6pt}
\smallskip
{\footnotesize\textcolor{celdark}{Part of the multimodal mechanics series
(analytical \textperiodcentered\ simulated \textperiodcentered\ experimental)
\,\textperiodcentered\, AIM Group, Imperial College London.}}

\end{document}
